\subsection{An alternative cycle-moment proof by entropy windows}
\label{n:coherent-window-route}

We give a second proof of the cycle-moment estimate
\eqref{n:all-cycle-moment}, retaining the node comparison and color
imbalance estimate already established.
Instead of separating dyadic bands into two parities, we control
their expectation under every tilted law of the output coins.  The
inverse-height weights in the entropy chain rule give each layer a
bounded total coefficient.  We retain the fixed input layers and the node cycle counts
$c_v$ from the proof of Theorem~\ref{n:strong-density}.

We first give a certificate version of the cycle estimate.  Consider
a butterfly of $q$ distinct-coordinate layers, with arbitrary fixed
permutations before and after it, and a fixed set $D$ of at most $k$
positions.  Let $Y$ count the cycles of the resulting permutation
contained in $D$, and let $r\ge1$ be an integer.
Put $\varepsilon=1/1000$.  If at least $r$ cycles lie in
$D$, with $r>k2^{-\varepsilon q}$, at least $r/2$ of them have
length at most $2k/r$.  A uniformly chosen closing path of one of
these cycles has at most $\log_2(2k/r)\le1+\varepsilon q$
internal contacts on average, by the contact bound used in
Section~\ref{n:cycle-encoding}.

Choose one root uniformly in each short cycle, and retain each
short cycle independently with probability $1/16$.  Conditional on
a particular cycle being retained, its closing path has at most
$q/16$ expected contacts with paths in other retained cycles:
at each layer there is only one other path at its switch, and its
cycle is retained with probability at most $1/16$.  Markov's
inequality shows that the probability of more than $q/4$ such
external contacts is at most $1/4$.  The probability of more than
$q/4$ internal contacts is at most $4/q+4\varepsilon$.
For all sufficiently large absolute $q$, the expected number of
retained roots passing both tests is therefore larger than
\[
 \frac1{16}\frac r2
          \left(1-\frac14-\frac4q-4\varepsilon\right)
 >\frac r{64}.
\]
Consequently some deterministic choice contains
$R=\lceil r/64\rceil$ cycles whose closing paths share at most
$q/2$ switches with all other paths in those selected cycles.
Deleting further selected cycles only improves this property.

A certificate specifies the unordered root set and the length of
each selected cycle; traverse the cycles in increasing root order.
There are at most
\[
 \binom kR(2k/r)^R
 \le2^{R\{C_0+2\varepsilon q\}}
\]
certificates.  Fix one certificate and expose its routes in that
order.  Before a closing route is queried, its initial point and
target are known.  The unique geometrical route between them is
therefore known too.  Whether it has at most $q/2$ already exposed
switches is measurable at this time.  If that test succeeds, at
least $q/2$ fresh fair bits must take prescribed values for the
cycle to close.  Its conditional cost is at most $2^{-q/2}$.
Iterating over the $R$ cycles and then summing certificates proves
\begin{equation}\label{n:window-certificate-tail}
 \Pr\{Y\ge r\}\le e^{-cqr}
 \quad\text{if }r>k2^{-\varepsilon q}
 \text{ and }q\ge q_0,
\end{equation}
for absolute $c>0$ and $q_0$.  This reasoning conditions only on
already exposed bits; it does not condition on a subsequently
discovered cycle.

Let $b_t=\sum_{v:\,\operatorname{height}(v)=t}c_v$ be the cycle
count at height $t$ in the there-and-back tree, and put
$q=\lfloor t/2\rfloor$.  Condition on all input coins and on output
coins at heights below $t-q$.  First suppose that $q\ge q_0$.
At a height-$t$ node, write the two child permutations as
$\alpha,\beta$, as in the cycle identification above.  Freeze
$\beta$ and everything in $\alpha$ except its last $q$ output
layers.  These are fresh distinct-coordinate layers,
so \eqref{n:window-certificate-tail} applies.  Integrating its tail,
and decreasing a positive absolute $a_1$ if necessary, gives
\[
 \log\mathbb E\left[e^{a_1t c_v}\mid
       \text{inputs},\text{output heights}<t-q\right]
       \le Ck_v e^{-a_2t}.
\]
The factor $t$ multiplying the exponentially small threshold is
absorbed by decreasing $a_2>0$.  Small heights are covered by
$c_v\le k_v$ and increasing $C$.  Nodes at one height use disjoint
remaining subtrees and $\sum_vk_v=l$.  Thus
\begin{equation}\label{n:coherent-window-mgf}
 \log\mathbb E\left[e^{a_1t b_t}\mid
       \text{inputs},\text{output heights}<t-q\right]
       \le Cl e^{-a_2t}.
\end{equation}

To make the change of law precise, fix the input coins and let $P$
be the product law of the output coins.  Order those coins by
increasing height.  For any probability $Q$ on this finite coin
space, put
\[
\begin{aligned}
 H_s&=\mathbb E_Q D\bigl(Q(\text{height-}s\text{ coins}\mid
           \text{lower heights})\,
          \|\,P(\text{height-}s\text{ coins})\bigr),\\
 H&=D(Q\|P)=\sum_sH_s.
\end{aligned}
\]
Here $D(\cdot\|\cdot)$ denotes relative entropy.  Applying the
entropy inequality conditionally to the window in
\eqref{n:coherent-window-mgf} gives, for $t\ge2$,
\[
 \mathbb E_Qb_t\le
 \frac{Cl e^{-a_2t}+
                  \sum_{s=t-\lfloor t/2\rfloor}^{t-1}H_s}{a_1t}.
\]
Here $b_t$ is measurable in output heights strictly below $t$.
The tilted law used in the conditional inequality is the marginal
law of the window given its lower-height history, with all future
coins integrated out.  The entropy chain rule says that its
conditional relative entropy, averaged over that history under $Q$,
is exactly the displayed sum of $H_s$.
The height-one term is at most $l$ directly.  A fixed layer $s$
belongs only to windows with $s+1\le t\le2s$, and
$\sum_{t=s+1}^{2s}t^{-1}\le\log2$.  Summation therefore yields
\begin{equation}\label{n:coherent-tilted-cycle-sum}
 \mathbb E_Q\sum_tb_t\le C_1l+C_2D(Q\|P).
\end{equation}
For $0<\delta\le1/C_2$, the finite-space entropy variational
formula now gives
\[
 \log\mathbb E_Pe^{\delta\sum_tb_t}
 =\sup_Q\left\{\delta\mathbb E_Q\sum_tb_t-D(Q\|P)\right\}
 \le\delta C_1l.
\]
The estimate is uniform in the fixed inputs, so it also holds after
averaging them.  Inserting it into the exact ideal-node recursion
\eqref{n:ideal-recursion}, with
$\theta\log2\le\delta$, proves the density theorem by this
alternative route; the smoothing and lower-diagram transfers then
apply without change.
